\section{本章定位与边界}
本章是数据截止后更新附录。全书封面基准日与\textbf{价格口径}仍为 \textbf{2026-06-26}（行情冻结）。对已发布 2026 H1 业绩预告的覆盖标的，报告采用\textbf{置信度优先}原则：管理层业绩预告是最高置信度的前瞻盈利信号，已在第 8 章\textbf{直接计入 2026E 归母分母、EPS、内在价值锚和综合目标价}（见第 8.2 节"业绩预告已计入估值"框）。本章的作用是给出\textbf{冻结口径（2026Q1 年化）与预告口径的并列对照}，让读者看清预告把估值分母抬升了多少、为什么综合目标价与评级最终落在何处，而不是把预告藏在附录里不计入。价格与未发预告标的的财务分母仍冻结在 2026-06-26。

\section{光通信板块 H1 2026 预告普查}
截至 2026-07-06，本报告 universe 共 36 只标的（估值覆盖 26 + 观察池 10），其中已发布正式 H1 2026 业绩预告的有 \textbf{3 只}。方法上，光通信概念板块接口在本机环境不可用；采用 akshare stock\_yjyg\_em(date=20260630) 业绩预告库 × 92 名光通信精选 universe 交叉核对，确认报告 36 只标的中的已披露名单。龙头（中际旭创、新易盛、天孚通信、光迅科技、源杰科技）尚未发布预告，其余标的完整半年报预约披露集中在 2026-08-01 至 08-31。其中永鼎股份（600105）与杭电股份（603618）为估值覆盖池标的，其 H1 预告已计入第 8 章估值分母；锐捷网络（301165）为观察池标的。

\begin{exhibitbox}[表：报告 universe 内已披露 H1 预告标的]
\centering
\small
\begin{tabularx}{\exhibitboxwidth}{>{\bfseries\raggedright\arraybackslash}p{1.6cm} >{\raggedright\arraybackslash}p{2.0cm} >{\centering\arraybackslash}p{1.6cm} >{\centering\arraybackslash}p{1.2cm} >{\raggedleft\arraybackslash}p{2.6cm} >{\raggedright\arraybackslash\sloppy\hspace{0pt}}X}
\toprule
\textbf{代码} & \textbf{名称} & \textbf{覆盖} & \textbf{类型} & \textbf{同比区间} & \textbf{公告日} \\
\midrule
600105 & 永鼎股份 & 估值覆盖 & 预增 & +57\%\textasciitilde{}+120\% & 2026-07-06 \\
603618 & 杭电股份 & 估值覆盖 & 预增 & +852\%\textasciitilde{}+958\% & 2026-07-04 \\
301165 & 锐捷网络 & 观察池 & 预增 & +32.7\%\textasciitilde{}+65.9\% & 2026-07-02 \\
\bottomrule
\end{tabularx}
\end{exhibitbox}

\section{H1 2026 预告明细}
下表汇总本轮已核实的 H1 2026 业绩预告，含 universe 外读数标的。归母净利单位为亿元，同比为公告披露区间。数据来源为公司公告与 akshare stock\_yjyg\_em 预告库交叉核对。

\begin{exhibitbox}[表：H1 2026 业绩预告明细]
\centering
\small
\begin{tabularx}{\exhibitboxwidth}{>{\bfseries\raggedright\arraybackslash}p{1.5cm} >{\raggedright\arraybackslash}p{1.8cm} >{\centering\arraybackslash}p{1.8cm} >{\raggedleft\arraybackslash}p{2.0cm} >{\raggedleft\arraybackslash}p{2.4cm} >{\centering\arraybackslash}p{1.7cm} >{\centering\arraybackslash}p{0.9cm}}
\toprule
\textbf{代码} & \textbf{名称} & \textbf{覆盖} & \textbf{H1归母(亿)} & \textbf{同比} & \textbf{公告日} & \textbf{来源} \\
\midrule
600105 & 永鼎股份 & 估值覆盖 & 5.00--7.00 & +57\%--+120\% & 2026-07-06 & S-35 \\
301165 & 锐捷网络 & 观察池 & 6.00--7.50 & +33\%--+66\% & 2026-07-02 & S-36 \\
603618 & 杭电股份 & 估值覆盖 & 3.60--4.00 & +852\%--+958\% & 2026-07-04 & S-37 \\
\bottomrule
\end{tabularx}
\sourcenote{来源 S-35/S-36/S-37 为公司业绩预告公告；M-04 为普查方法记录。永鼎 600105 与杭电 603618 为估值覆盖标的，H1 预告已计入第 8 章估值分母；锐捷 301165 为观察池标的。}
\end{exhibitbox}


\section{永鼎股份（600105）：业绩预告已计入估值（冻结口径对照）}
永鼎股份 属估值覆盖池，2026-07-06 披露 H1 2026 归母净利 5.00--7.00 亿元、同比 +57\%--+120\%（预增）。该预告已按置信度优先原则\textbf{计入第 8 章估值分母与目标价}；下表给出采用口径与冻结（Q1 年化）口径的并列对照，说明预告把估值分母抬升了多少。H1 预告中值 6.00 亿元按 H1 占全年 0.50 年化得 2026E 归母约 12.00 亿元，并按量价齐升在 Q1 净利率下同步上抬 2026E 收入（Q1 年化归母口径为 6.62 亿元）；Q1 归母同比 -45.2\%（承压）到 H1 同比 +57\%\textasciitilde{}+120\%（大幅预增），确认 Q2 强反转。

\begin{exhibitbox}[估值卡：永鼎股份 冻结 vs 采用]
\centering
\small
\begin{tabularx}{\exhibitboxwidth}{>{\bfseries\raggedright\arraybackslash}p{3.2cm} >{\raggedleft\arraybackslash}p{2.7cm} >{\raggedleft\arraybackslash}p{3.3cm} >{\raggedright\arraybackslash\sloppy\hspace{0pt}}X}
\toprule
\textbf{项目} & \textbf{冻结口径 (Q1 年化)} & \textbf{采用口径 (H1 预告，已入估值)} & \textbf{说明} \\
\midrule
2026E 归母净利(亿) & 6.62 & 12.00 & H1 中值按 H1 占全年 0.50 年化，作为估值分母 \\
2026E 收入(亿) & 51.90 & 94.12 & 量价齐升，收入按 Q1 净利率与指引利润同步上抬（保守下限） \\
2026E EPS & 0.45 & 0.82 & EPS 近乎翻倍，反映 Q2 强反转 \\
收入/股 & 3.55 & 6.44 & PS 腿（权重 50\%）随收入上修，非仅利润修正 \\
2027E EPS & 0.50 & 0.90 & 按公司层面成长假设外推 \\
内在价值锚 & 14.77 & 24.57 & PE/PB/PS 各腿均随预告上修，内在锚显著抬升 \\
综合目标价 & 46.00 & 46.00 & 现价仍较内在锚溢价约 +167\%，市场情绪锚地板支撑，目标价基本持平 \\
评级 & 中性观察 & 中性观察 & 评级维持不变 \\
盈利质量标签 & 承压 & Q2 反转确认 & Q1 同比 -45.2\% 承压 → H1 大幅预增 \\
\bottomrule
\end{tabularx}
\end{exhibitbox}

预告已计入 EPS 与 PE/PB/PS 各条估值腿，内在价值锚显著上修、估值泡沫度收窄；但评级不因单纯利润超预期机械上修：永鼎当前价即使在预告口径下仍较内在锚溢价约 +167\%，受市场情绪锚地板支撑，综合目标价与中性观察维持不变，需 Q2/Q3 收入、毛利率和现金流同步确认后才具备评级上修条件。 因此本次采用预告口径为\textbf{盈利质量与 EPS 层面的上修}，而非评级上修：待 Q2/Q3 收入、毛利率与现金流同步确认后，方具备将 永鼎股份 从中性观察上修的条件。组合权重维持不变，价格口径仍为 2026-06-26。


\section{杭电股份（603618）：业绩预告已计入估值（冻结口径对照）}
杭电股份 属估值覆盖池，2026-07-04 披露 H1 2026 归母净利 3.60--4.00 亿元、同比 +852\%--+958\%（预增）。该预告已按置信度优先原则\textbf{计入第 8 章估值分母与目标价}；下表给出采用口径与冻结（Q1 年化）口径的并列对照，说明预告把估值分母抬升了多少。H1 预告中值 3.80 亿元按 H1 占全年 0.50 年化得 2026E 归母约 7.60 亿元，并按量价齐升在 Q1 净利率下同步上抬 2026E 收入（Q1 年化归母口径为 3.37 亿元）；Q1 归母同比 +280.0\%（承压）到 H1 同比 +57\%\textasciitilde{}+120\%（大幅预增），确认 Q2 强反转。

\begin{exhibitbox}[估值卡：杭电股份 冻结 vs 采用]
\centering
\small
\begin{tabularx}{\exhibitboxwidth}{>{\bfseries\raggedright\arraybackslash}p{3.2cm} >{\raggedleft\arraybackslash}p{2.7cm} >{\raggedleft\arraybackslash}p{3.3cm} >{\raggedright\arraybackslash\sloppy\hspace{0pt}}X}
\toprule
\textbf{项目} & \textbf{冻结口径 (Q1 年化)} & \textbf{采用口径 (H1 预告，已入估值)} & \textbf{说明} \\
\midrule
2026E 归母净利(亿) & 3.37 & 7.60 & H1 中值按 H1 占全年 0.50 年化，作为估值分母 \\
2026E 收入(亿) & 91.77 & 207.05 & 量价齐升，收入按 Q1 净利率与指引利润同步上抬（保守下限） \\
2026E EPS & 0.50 & 1.13 & EPS 近乎翻倍，反映 Q2 强反转 \\
收入/股 & 13.62 & 30.73 & PS 腿（权重 50\%）随收入上修，非仅利润修正 \\
2027E EPS & 0.57 & 1.30 & 按公司层面成长假设外推 \\
内在价值锚 & 12.65 & 22.03 & PE/PB/PS 各腿均随预告上修，内在锚显著抬升 \\
综合目标价 & 20.29 & 25.80 & 现价仍较内在锚溢价约 +167\%，市场情绪锚地板支撑，目标价基本持平 \\
评级 & 减持 & 减持 & 评级维持不变 \\
盈利质量标签 & 承压 & Q2 反转确认 & Q1 同比 +280.0\% 承压 → H1 大幅预增 \\
\bottomrule
\end{tabularx}
\end{exhibitbox}

预告已计入 EPS 与 PE/PB/PS 各条估值腿，内在价值锚显著上修、估值泡沫度收窄；但评级不因单纯利润超预期机械上修：永鼎当前价即使在预告口径下仍较内在锚溢价约 +167\%，受市场情绪锚地板支撑，综合目标价与中性观察维持不变，需 Q2/Q3 收入、毛利率和现金流同步确认后才具备评级上修条件。 因此本次采用预告口径为\textbf{盈利质量与 EPS 层面的上修}，而非评级上修：待 Q2/Q3 收入、毛利率与现金流同步确认后，方具备将 杭电股份 从中性观察上修的条件。组合权重维持不变，价格口径仍为 2026-06-26。
